% Source: https://artofproblemsolving.com/community/c3h2198006_probability_question
% Stable identity: aops-topic-2198006-q01
The question can only be consistent if the declared value equals the probability of picking a choice that shows this value.
\begin{itemize}
\item If $25\%$ is correct, then (A) and (D) are both correct choices, so the chance of picking a correct choice is $\dfrac{2}{4}=50\%\ne25\%$.
\item If $50\%$ is correct, only (C) is correct, so the chance is $\dfrac{1}{4}=25\%\ne50\%$.
\item If $0\%$ is correct, then (B) is a correct choice, so the chance is $\dfrac{1}{4}=25\%\ne0\%$.
\item If no choice is correct, the chance of being correct is $0\%$; but then (B) states the true value and is correct after all, a contradiction.
\end{itemize}
Every possibility contradicts itself, so the question has no consistent answer. It is a self-referential paradox, like ``this sentence is false''.
